\section{Arquitectura de referencia}

La arquitectura de referencia separa la captura de ubicación, la autenticación, el almacenamiento y la visualización. La Figura~\ref{fig:arquitectura} es un espacio para que el equipo incluya su diagrama; la Tabla~\ref{tab:componentes-arquitectura} describe las responsabilidades de cada componente.

\begin{figure}[H]
	\centering
	\fbox{%
		\parbox[c][6.2cm][c]{0.92\textwidth}{%
			\centering
			\textbf{[Insertar aquí el diagrama de arquitectura de la solución]}\\[0.4cm]
			\footnotesize
			Incluya la aplicación móvil, el bot de Telegram, la API REST, el servicio de autenticación,\\
			la base de datos, el servicio de mapas y el panel administrativo.\\
			Señale con flechas el sentido del flujo y los controles TLS en cada tramo.
		}%
	}
	\caption{Arquitectura de referencia de la solución}
	\label{fig:arquitectura}
\end{figure}

\begin{table}[H]
	\caption{Componentes de la arquitectura y responsabilidades}
	\label{tab:componentes-arquitectura}
	\centering
	\footnotesize
	\renewcommand{\arraystretch}{1.2}
	\begin{tabularx}{\textwidth}{@{}>{\raggedright\arraybackslash}p{0.17\textwidth} Y >{\raggedright\arraybackslash}p{0.16\textwidth} >{\raggedright\arraybackslash}p{0.15\textwidth} >{\raggedright\arraybackslash}p{0.15\textwidth}@{}}
		\toprule
		\textbf{Componente} & \textbf{Función} & \textbf{Entrada} & \textbf{Salida} & \textbf{Tecnología} \\
		\midrule
		Aplicación móvil & Captura y envía la ubicación con consentimiento. & GPS, usuario & Coordenadas, estado & \\
		Bot de Telegram & Solicita la ubicación de forma manual. & Comando del usuario & Coordenadas & \\
		API REST & Valida, autentica y registra las solicitudes. & Peticiones HTTPS & Respuestas JSON & \\
		Autenticación & Verifica identidad y permisos. & Número, token & Sesión / JWT & \\
		Base de datos & Persiste usuarios, dispositivos, ubicaciones y consentimientos. & Datos validados & Consultas & \\
		Servicio de mapas & Representa las ubicaciones y resuelve direcciones. & Coordenadas & Mapa / dirección & \\
		Panel administrativo & Consulta y audita el sistema. & Credenciales & Reportes & \\
		\bottomrule
	\end{tabularx}
\end{table}

\section{Modelo de datos propuesto}

La Tabla~\ref{tab:modelo-datos} propone las entidades mínimas. El equipo debe ajustar los campos a su implementación y documentar las claves primarias, foráneas e índices.

\begin{table}[H]
	\caption{Entidades y campos principales del modelo de datos}
	\label{tab:modelo-datos}
	\centering
	\footnotesize
	\renewcommand{\arraystretch}{1.2}
	\begin{tabularx}{\textwidth}{@{}>{\raggedright\arraybackslash}p{0.17\textwidth} Y >{\raggedright\arraybackslash}p{0.25\textwidth}@{}}
		\toprule
		\textbf{Entidad} & \textbf{Campos principales} & \textbf{Relaciones} \\
		\midrule
		Usuario & Identificador, nombre, número telefónico, correo, hash de contraseña, estado. & Uno a muchos con dispositivos y consentimientos. \\
		Dispositivo & Identificador, usuario, número, plataforma, token, estado. & Muchos a uno con usuario; uno a muchos con ubicaciones. \\
		Ubicación & Identificador, dispositivo, latitud, longitud, precisión, fecha/hora, fuente. & Muchos a uno con dispositivo. \\
		Consentimiento & Identificador, usuario, finalidad, alcance, fecha, estado, fecha de revocación. & Muchos a uno con usuario. \\
		Zona & Identificador, usuario, nombre, centro, radio. & Uno a muchos con alertas. \\
		Alerta & Identificador, zona, dispositivo, tipo, fecha, estado. & Muchos a uno con zona y dispositivo. \\
		& & \\
		\bottomrule
	\end{tabularx}
\end{table}

\section{Plan de trabajo}

El plan de la Tabla~\ref{tab:plan-trabajo} distribuye las tres semanas del taller. Ajuste las fechas y los responsables según la disponibilidad del equipo; conserve los hitos de revisión y la entrega final.

\begin{table}[H]
	\caption{Plan de trabajo sugerido}
	\label{tab:plan-trabajo}
	\centering
	\footnotesize
	\setlength{\tabcolsep}{3pt}
	\renewcommand{\arraystretch}{1.2}
	\begin{tabularx}{\textwidth}{@{}c >{\raggedright\arraybackslash}p{0.12\textwidth} Y >{\raggedright\arraybackslash}p{0.20\textwidth} >{\raggedright\arraybackslash}p{0.12\textwidth}@{}}
		\toprule
		\textbf{Semana} & \textbf{Fechas} & \textbf{Actividades} & \textbf{Entregable} & \textbf{Responsable} \\
		\midrule
		1 & 26 oct -- 1 nov & Marco del laboratorio, topología, casos 1 a 3 y diseño inicial de la Parte~II. & Fichas 1--3 y propuesta inicial. & \\
		2 & 2 -- 8 nov & Casos 4 a 7, arquitectura, modelo de datos y API. & Fichas 4--7 y arquitectura. & \\
		3 & 9 -- 14 nov & Casos 8 a 10, implementación, pruebas, informe final y hashes. & Versión final y anexos. & \\
		\bottomrule
	\end{tabularx}
\end{table}

\section{Entregables de la Parte II}

\begin{itemize}[leftmargin=*, topsep=2pt, itemsep=4pt, label=\checkbox]
	\item Código fuente funcional en un repositorio con historial de commits.
	\item README con requisitos, instalación, configuración, ejecución y pruebas.
	\item Diagrama de arquitectura y modelo de datos.
	\item Documentación de la API REST (OpenAPI o equivalente).
	\item Demostración con datos sintéticos o de laboratorio (video, capturas o ejecución en vivo).
	\item Plan de pruebas y resultados de los casos felices y de error.
	\item Política de privacidad, consentimiento informado y mecanismo de revocación.
	\item Análisis de riesgos, limitaciones y trabajo futuro.
\end{itemize}
